\documentclass{beamer}

%-----------------------------------------------
%Pacotes: Insira aqui os pacotes a serem usados 

\usepackage[utf8]{inputenc}
\usepackage[brazil]{babel}
\usepackage{amssymb, amsmath, amsthm, amsfonts, amscd, graphicx, subfigure, array, multicol, setspace, caption, color}
\newtheorem{Teorema}{Teorema}
\newtheorem{Lema}{Lema}
\newtheorem{Corolário}{Corolário }
\newtheorem{Proposição}{Proposição}

 %----------------------------------------------
 %Comandos da capa: Adaptar com suas informações

\title[Título]{Título}
\author[Nome do Autor]{Nome do autor\\ Orientador: Nome do Orientador \\Coorientador: Nome do Coorientador}
\institute[Instituição]{II Mostra IC de Matemática \\ Universidade Federal de Ouro Preto}
\date[II Mostra IC de Matemática]{2024}

 %----------------------------------------------
 %Configurações do template: NÃO EDITAR
 
\usetheme{Madrid}
\definecolor{bege}{rgb}{0.88, 0.75, 0.54}
\setbeamercolor{title}{fg=white, bg=black}
\setbeamercolor{title in head/foot}{fg=black,bg=bege}
\setbeamercolor{author in head/foot}{fg=white,bg=black}
\setbeamercolor{frametitle}{fg=white,bg=black}
\setbeamercolor{block title}{fg=white,bg=black}
\setbeamercolor{block body}{bg=bege, fg=black}
\setbeamercolor{date in head/foot}{fg=white, bg=black}
\setbeamercolor{item}{fg=bege}
\setbeamertemplate{navigation symbols}{}

%-----------------------------------------------
%Página de capa: NÃO EDITAR O CÓDIGO ABAIXO (as devidas alterações já devem ter sido feitas nos comandos de capa)

\begin{document}
\begin{frame}
    \begin{picture}(3.5cm,3.5cm) 
        \put(10,40){\includegraphics[scale=0.15]{Logos/Logo UFOP.png}}
        \put(137,20){\includegraphics[scale=0.09]{Logos/Logo Mostra IC.png}}
        \put(300,50){\includegraphics[scale=0.05]{Logos/Logo Demat.png}}
        \put(300,30){\includegraphics[scale=0.16]{Logos/Logo PET.png}}
    \end{picture}
    \maketitle
\end{frame}

 %---------------------------------------------
 %Apresentação: Totalmente livre para edição

\begin{frame}{Título do Frame}

    \begin{block}{Exemplo de bloco}
    Este é um exemplo de bloco.
    \end{block}
    
    \begin{figure}
        \centering
        \includegraphics[scale=0.1]{Logos/Logo Demat.png}
        \caption{Exemplo de figura}
        \label{ex.figura}
    \end{figure}
    
\end{frame}

\begin{frame}{Título do Frame}

    \begin{itemize}
        \item Este
        \item é
        \item um
        \item exemplo
        \item de
        \item lista
    \end{itemize}
    
    Para citações no texto, usar: \cite{burgos2016solution}, \cite{de2001edp}, \cite{Leibniz}, \cite{Eder},
    
\end{frame}

\begin{frame}{Título do Frame}
    \begin{Teorema}
        Enunciado do Teorema.
    \end{Teorema}

    \begin{Corolário}
        Escreva também corolários, proposições e lemas.

    \end{Corolário}
    
\end{frame}

%----------------------------------------------
%Bibliografia: Manter o código abaixo sem alterações. Para acrescentar e/ou remover referências, editar o arquivo "ref.bib".

\begin{frame}[allowframebreaks]{Bibliografia}
    \bibliographystyle{plain}
    \bibliography{ref.bib}
\end{frame}

\end{document}